\title{Parameter-Efficient Orthogonal Finetuning via Factorization}

\begin{abstract}
The prevalence of large foundation models has led to significant advancements, but their immense computational demands make training from scratch impractical. As such, methods to adapt these models efficiently for specific downstream tasks have become increasingly valuable. In this work, we investigate Orthogonal Finetuning (OFT) as a systematic approach for model adaptation. Although OFT offers strong generalization performance, it still necessitates training a considerable number of parameters due to the large size of orthogonal matrices. To mitigate this, we analyze OFT through the lens of information transmission, uncovering essential criteria that underpin greater parameter efficiency. Drawing inspiration from the Cooley-Tukey fast Fourier transform, which enables efficient data transformation, we propose a novel parameterization of orthogonal matrices utilizing butterfly structures. This new parameterization is integrated into OFT, resulting in a new approach we call Orthogonal Butterfly (BOFT), which achieves parameter-efficient finetuning. BOFT generalizes OFT, encompassing it as a specific instance, and thereby yields a more flexible framework for orthogonal finetuning. We then perform comprehensive experiments, adapting large vision transformers, large language models, and text-to-image diffusion architectures to a diverse range of vision and language tasks.
\end{abstract}

\section{Introduction}

Cutting-edge systems such as ChatGPT~\cite{brown2020language,chen2021evaluating} and Stable Diffusion~\cite{rombach2022high} exemplify the robust generalization capabilities of contemporary large-scale foundation models. However, the extraordinary rise in their effectiveness has been matched by a substantial escalation in the number of parameters

(\emph{e.g.}, GPT-3 contains approximately 175B parameters), rendering the process of training these models from the beginning an overwhelming endeavor for most researchers. Consequently, advancing the field relies heavily on methods that allow for effective adaptation of existing models, bypassing the need for retraining from scratch. Thus, efficient transfer of foundation models to downstream tasks is essential. Current strategies for such adaptation can be grouped into three primary categories: \emph{model finetuning}~\cite{hulora2022,zaken2022bitfit,guo2020parameter,raffel2020exploring,qiu2023controlling,zhang2022adaptive,chavan2023one}, where selective parameters within the model are updated while the underlying architecture stays the same; \emph{adapter tuning}~\cite{rebuffi2017learning,houlsby2019parameter,pfeiffer2020adapterfusion,lin2020exploring,he2021towards}, which introduces supplementary trainable components into the model; and \emph{prompt tuning}~\cite{li2021prefix,lester2021power}, which augments the model with tunable prefix tokens at the input level. Among these, model finetuning is particularly notable for its simplicity and effectiveness, as well as for not incurring any additional latency during inference.

The core idea underlying model finetuning is to maintain similarity between the pretrained model and its finetuned counterpart according to certain criteria, thereby safeguarding the acquired knowledge from pretraining. A typical practice in current finetuning methodologies is the use of a small learning rate, an intuitive strategy aimed at limiting the divergence between the weights of the pretrained and finetuned models. Given that weight matrices possess an inherent structure, a more systematic method focuses on conserving relational properties of the weights, specifically the pairwise angular relationships between neurons. Building on this observation, the Orthogonal Finetuning~(OFT) framework~\cite{qiu2023controlling} was proposed. In OFT, all neurons within the same layer undergo a transformation via a common orthogonal matrix, guaranteeing that their pairwise angular relationships remain invariant throughout finetuning. While OFT achieves notable gains in generalization and convergence for text-to-image diffusion model finetuning~\cite{qiu2023controlling}, a significant limitation is the extensive parameter count incurred by the large dimensions of the required orthogonal matrices. To mitigate this, OFT employs a block-diagonal decomposition, reducing the number of trainable parameters. Nevertheless, this parameter efficiency is achieved by enforcing a rigid sparsity pattern on the orthogonal matrix, applying transformations independently within separate blocks. Such a block-diagonal structure, although economizing on parameters, can inject unwarranted inductive biases (e.g., decreased representational capacity and the inability to recover traditional linear transformations), and, importantly, the challenge of selecting an optimal sparsity configuration remains unresolved.

A central challenge here lies in constructing a parameter-efficient dense orthogonal matrix. Although, at first glance, this might appear unattainable due to the requirement of $\mathcal{O}(d^2)$ parameters for a $d$-dimensional dense orthogonal matrix, our method introduces a distinctive strategy: we assemble a dense orthogonal matrix by sequentially multiplying several factorized sparse matrices. This method is founded on the principle that computational resources (increased multiplication steps per training iteration) can be substituted for storage space (fewer parameters). Since our orthogonal matrix is represented as the product of sparse matrices, we must execute these matrix multiplications repeatedly during every training cycle. To frame this process, we conceptualize the formation of a dense orthogonal matrix as an information propagation problem, wherein the matrix generation process via sparse matrix multiplications is akin to transmitting signals across a graph with a grid topology. Such a perspective allows us to devise multiple strategies for sparse matrix factorization, restricting the number of trainable parameters while retaining the capacity to express dense matrices effectively.

To further economize on parameters within our information transmission paradigm, we turn to the butterfly patterns found in the Cooley-Tukey fast Fourier transform algorithm~\cite{cooley1965algorithm}, which demonstrate efficient inter-node communication~\cite{klasing1994broadcasting}. Specifically, the butterfly graph topology inherent to the Cooley-Tukey scheme leads naturally to the "butterfly factorization" method for sparse matrices. Utilizing butterfly factorization, it becomes feasible to synthesize a $d$-dimensional dense matrix by multiplying $\mathcal{O}(\log d)$ sparse matrices, each featuring only $\mathcal{O}(d)$ non-zero elements. By maintaining the orthogonality of each sparse factor, we are able to achieve a compact orthogonal matrix representation with just $\mathcal{O}(d\log d)$ parameters—marking a significant improvement in parameter efficiency. Building upon this representation, we introduce Orthogonal Butterfly~(BOFT), a new approach for parameter-efficient finetuning. BOFT encompasses OFT as a particular instance and offers a comprehensive framework for orthogonal finetuning. Notably, both the block-diagonal and butterfly structures share the property of sparsity, utilizing sparse arrangements to reduce the number of tunable parameters in orthogonal matrices. Comparing our approach to the low-rank strategies used in LoRA~\cite{hulora2022} is also insightful, as both low-rank and sparse matrices are prominent structured matrix families~\cite{candes2011robust} known for facilitating parameter-efficient designs.

In contrast to OFT's reliance on a block-diagonal structure, which negotiates between expressive power and regularity, BOFT leverages the butterfly structure to generate a smoother continuum between fully expressive matrices from the complete orthogonal group and regularizing identity matrices. This strategy allows us to identify a more compact hypothesis space within the orthogonal group suited to downstream tasks. Given the prominent role of the butterfly structure in expediting many classic linear operations—including the discrete Fourier and discrete cosine transforms—we posit that our adoption of such structural constraints for parameter efficiency inherently instills a beneficial inductive bias, fostering improved generalization and reducing the likelihood of overfitting. Our rationale is rooted in the butterfly structure’s capacity to readily reproduce many standard linear transforms, a capability not shared by the block-diagonal architecture employed in OFT. The core contributions of this work are as follows:

\begin{itemize}[leftmargin=*,nosep]
    \item We analyze the challenge of parameter efficiency in orthogonal finetuning by introducing an innovative information transmission paradigm. This approach reframes the search for a parameter-efficient dense orthogonal matrix as an information propagation task in a grid-structured network.
    \item Drawing inspiration from the butterfly configurations utilized in the Cooley-Tukey algorithm, we introduce Orthogonal Butterfly—a novel, parameter-efficient method for orthogonal finetuning—developed within the information transmission perspective.
    \item We furnish theoretical arguments elucidating how BOFT achieves substantial reductions in trainable parameters without sacrificing expressiveness or training stability. The matrix factorization inherent to BOFT further enables a compelling form of weight interpolation (see Figure~\ref{fig:boft_interpolation}).
    \item For the first time, we extend orthogonal finetuning to a wide array of tasks beyond controllable text-to-image generation~\cite{qiu2023controlling}, demonstrating its strong versatility as a general-purpose finetuning strategy. Specifically, BOFT is deployed across diverse adaptation scenarios in both computer vision and natural language processing, consistently outperforming leading contemporary methods in terms of parameter efficiency and generalization performance.
\end{itemize}

\section{Related Work}

\textbf{Parameter-efficient finetuning (PEFT)}. As the scale and capabilities of foundation models (\emph{e.g.}, \cite{brown2020language,radford2021learning,rombach2022high,kirillov2023segment}) continue to expand, there has been a surge of research interest in developing approaches for finetuning these models on downstream tasks with minimal parameter updates~\cite{treviso2023efficient,ding2023parameter,lialin2023scaling}. A wide array of PEFT strategies have emerged~\cite{houlsby2019parameter,aghajanyan2020intrinsic,hulora2022,edalati2022krona,wang2022adamix,gheini2021cross,zaken2022bitfit,guo2020parameter,sung2021training,ansell2022composable,lester2021power,li2021prefix,vu2022spot,he2021towards,mao2021unipelt,karimi2021compacter,liu2022few,sung2022lst,chen2022parameter,jia2022visual,chen2022adaptformer,zhang2022neural,jie2023fact,lian2022scaling,luo2023towards}; of these, approaches utilizing reparameterization~\cite{aghajanyan2020intrinsic,hulora2022,edalati2022krona} are especially pertinent to our study. LoRA~\cite{hulora2022}, for example, enhances the pretrained weight matrices by incorporating a low-rank adaptation term constructed from the multiplication of two small-rank matrices. This methodology has shown substantial success on NLP benchmarks. Since LoRA assigns a uniform rank to the inserted matrices across all layers, later efforts~\cite{zhang2022adaptive,valipour2022dylora,zhang2023increlora} have explored adaptive rank allocation to optimize the use of available parameters per layer. The efficiency and simplicity of this low-rank formulation have led to its widespread adoption \cite{dettmers2023qlora,zi2023delta,chavan2023one}. Inspired by studies on the influence of hyperspherical energy in understanding generalization~\cite{liu2018learning,liu2021orthogonal}, orthogonal finetuning (OFT) was proposed in \cite{qiu2023controlling} as a powerful alternative for text-to-image diffusion model adaptation. In OFT, the approach is to optimize an orthogonal transformation matrix that maps neurons within a layer, resulting in better generalization performance and more robust training than LoRA. Nevertheless, OFT typically involves a larger number of trainable parameters compared to LoRA, indicating a need for enhanced parameter efficiency. Furthermore, the broader applicability of OFT to tasks beyond text-to-image diffusion remains largely unexplored. BOFT addresses this efficiency gap in OFT by leveraging butterfly factorization, making it possible to apply orthogonal finetuning more broadly across different adaptation scenarios.

\textbf{Butterfly structures}. The radix-2 Cooley-Tukey algorithm efficiently decomposes the $N$-point discrete Fourier transform into recursive computations on two subproblems of size $\frac{N}{2}$, naturally yielding a butterfly structure. Such a structure can be mathematically represented as the product of multiple sparse matrices, collectively forming a butterfly matrix. Butterfly matrices have found applications as orthogonal parameterizations that circumvent pivoting during Gaussian elimination and enable computational gains~\cite{parker1995random}; they have also contributed to stable training in recurrent neural architectures~\cite{jing2017tunable} and efficient kernel approximation~\cite{munkhoeva2018quadrature}. Advances by \cite{dao2019learning,dao2019kaleidoscope} involve learning efficient linear transformations through butterfly matrix parametrizations, while subsequent work \cite{chen2022pixelated,dao2022monarch} employs butterfly matrices to achieve neural network training with reduced computational costs. Their utility extends further to fast matrix-vector multiplications~\cite{michielssen1996multilevel,o2010algorithm}, compact representations for large matrices~\cite{li2015butterfly}, and enhanced network transmission protocols~\cite{klasing1994broadcasting,li2003linear,rajasingh2016transmission}. Our research diverges from these previous uses by focusing on integrating butterfly structures to increase the parameter efficiency of OFT.

\section{An Information Transmission View on Orthogonal Finetuning}

We begin by reviewing the foundational aspects of OFT. In OFT, adaptation of pretrained weights is accomplished by expressing the updated weight matrix as the multiplication of a trainable orthogonal matrix with the fixed pretrained weight matrix. This approach stands in contrast to LoRA, which applies parameter updates through the addition of a low-rank matrix, while OFT employs the multiplicative influence of an orthogonal matrix instead. LoRA achieves parameter-efficiency via a low-rank decomposition, whereas standard OFT attains this through a block-diagonal configuration of the orthogonal matrix~\cite{qiu2023controlling}; the smaller the individual block size, the more economical the parameter usage in OFT. An illustrative comparison is provided in Figure~\ref{fig:comp_lora}. The rationale for using orthogonal transformations during finetuning is to maintain the angular relationships between neuron outputs~\cite{liu2018learning,liu2021orthogonal,qiu2023controlling}, thereby retaining much of the semantic information learned during pretraining. Specifically, OFT learns an orthogonal matrix $\bm{R}\in\mathbb{R}^{d\times d}$ for a pretrained linear transformation $\bm{W}^0\in\mathbb{R}^{d\times n}$, adjusting the forward computation from $\bm{z}=(\bm{W}^0)^\top\bm{x}$ to $\bm{z}=(\bm{R}\bm{W}^0)^\top\bm{x}$, where $\bm{x}\in\mathbb{R}^{d}$ is the input and $\bm{z}\in\mathbb{R}^{n}$ is the output. To allow for zero initialization, $\bm{R}$ is initially set to the identity matrix. To retain the orthogonal constraint on $\bm{R}$ during optimization, following \cite{qiu2023controlling,liu2021orthogonal}, Cayley parameterization is adopted: $\bm{R}=(\bm{I}+\bm{Q})(\bm{I}-\bm{Q})^{-1}$, where $\bm{Q}$ denotes a skew-symmetric matrix satisfying $\bm{Q}=-\bm{Q}^\top$. For enhanced parameter-efficiency, the block-diagonal architecture expresses $\bm{R}$ as $\text{diag}(\bm{R}_1,\bm{R}_2,\cdots,\bm{R}_r)$, with each $\bm{R}_i\in\mathcal{R}^{b\times b}$ being a small orthogonal block, and $br=d$. The parameter saving of this design stems from the premise that neuron dimensions (i.e., elements within a column of $\bm{W}^0$) are partitioned into $r$ segments, with each segment being transformed via an independent orthogonal matrix. Nonetheless, despite demonstrated empirical strengths, such block-wise division according to indices lacks inherent justification from a neuronal perspective, rendering dense orthogonal matrices a compelling alternative. This raises a key inquiry: \emph{Is it possible to design a dense orthogonal matrix structure while preserving parameter-efficiency?}

In order to tackle this issue, we suggest representing a dense orthogonal matrix as a product of several sparse orthogonal matrices. To offer a coherent and accessible framework for analyzing the sparsity patterns involved in orthogonal matrix decompositions, we recast the generation of a dense orthogonal matrix in OFT as an information transmission process. Concretely, the task of constructing a dense matrix $\bm{R}\in\mathbb{R}^{d\times d}$ using $m$ square matrices, so that $\thickmuskip=2mu \medmuskip=2mu \bm{R}=\bm{B}_m\bm{B}_{m-1}\cdots\bm{B}_1$, can be analogized to information passing through a network structured as a grid with $d\times (m+1)$ nodes, as depicted in Figure~\ref{fig:info_trans}. This information transmission perspective is motivated by the idea that a dense square matrix of dimension $d$ represents full connectivity between one set of $d$ nodes and another. In this interpretation, a zero in position $R_{ij}$ of $\bm{R}$ signals that no information is transmitted from node $j$ to node $i$, while a nonzero entry permits the flow of information. Thus, viewing $\bm{R}$ as the product $\bm{B}_m\bm{B}_{m-1}\cdots\bm{B}_1$ is equivalent to modeling a sequence of information exchanges, each defined by the graphs determined by $\bm{B}_i$ for all $i$. Information propagates layer by layer, starting from the connectivity dictated by $\bm{B}_1$ up to $\bm{B}_n$ at the last stage. As an illustrative scenario in Figure~\ref{fig:info_trans}, the factorization $\bm{R}=\bm{B}_5\bm{B}_4\bm{B}_3\bm{B}_2\bm{B}_1$ is shown alongside its associated sparsity patterns and the corresponding graph structure. This graphical representation arises by expanding the process of matrix multiplication. Each $\bm{B}_i$ defines how nodes at level $i$ are connected to nodes at level $i+1$, where the $(j_1,j_2)$ entry of $\bm{B}_i$ encodes the presence (nonzero) or absence (zero) of a directed connection from node $j_2$ in layer $i$ to node $j_1$ in layer $i+1$. For the composite matrix $\bm{B}_5\bm{B}_4\bm{B}_3\bm{B}_2\bm{B}_1$ to produce a dense result, it is essential that every source node in the first level can reach all target nodes in the sixth level via the sequence of matrix products. If we restrict our attention to $\bm{R}=\bm{B}_4\bm{B}_3\bm{B}_2\bm{B}_1$, linking nodes from level one to level five, it may be observed (for instance) that node $1$ cannot transmit to node $3$—meaning that the entry at $(3,1)$ in the corresponding product is zero. The same relationship between the resultant graph structure and the matrix's sparsity pattern applies for shorter products like $\bm{R}=\bm{B}_3\bm{B}_2\bm{B}_1$ and $\bm{R}=\bm{B}_2\bm{B}_1$. In a more general setting, for $\bm{R}\in\mathbb{R}^{d\times d}$ to be dense, the collection of $m$ factors $\bm{B}_m,\dots,\bm{B}_1$ must correspond to a set of directed edges on a $d\times (m+1)$ grid, where edges only connect nodes in neighboring columns (levels), so that information is able to flow from every node at the initial level to each node at the final $(m+1)$-th level.

Adopting the information transmission perspective offers deeper insight into the sparsity patterns found in the factorization matrices of OFT. In Figure~\ref{fig:blk_diag}, the block-diagonal layout of $\bm{R}$ in standard OFT is depicted. Although this structure results in fewer trainable parameters, it restricts $\bm{R}$ to a non-dense configuration. Our objective is to synthesize a dense orthogonal matrix by leveraging $m$ sparse orthogonal matrices, while utilizing the smallest number of effective trainable parameters. Guided by the information transmission view, our requirements are as follows: \emph{(i)} \textbf{dense connectivity}, ensuring each node at the first layer has at least one path reaching every node in the final layer, and \emph{(ii)} \textbf{minimal free edges}, where the sum of all edges is minimized, subject to the orthogonality constraint. It is important to recognize that orthogonality imposes specific constraints on edge connections between adjacent layers. For instance, to guarantee that each $\bm{B}_i$ is full-rank (essential for orthogonality), exactly $d$ edges are required to establish a bijective relationship between the nodes in level $i$ and those in level $(i-1)$; thus, at least $d$ edges must exist between each pair of consecutive layers (as illustrated, e.g., by 4 edges in Figure~\ref{fig:info_trans}). Since these $d$ edges are fixed by orthogonality and cannot be trained, they should be excluded from the count of trainable edges; for instance, a $d \times d$ orthogonal matrix containing $d$ non-zero entries only permits those elements to take values of $\pm 1$. As orthogonal matrices require fewer free parameters than dense matrices, imposing orthogonality introduces further dependencies among the edges. For example, in a $2 \times 2$ orthogonal matrix, setting a zero in position $(1,1)$ necessitates another zero in position $(2,2)$—thus, omitting a single edge can force the omission of another. Therefore, depending on a particular edge configuration, orthogonality may lead to the addition or removal of edges. Examining the pattern of non-zero entries in the constructed orthogonal matrix, the information transmission framework proves especially illuminating for OFT, as the primary concern is with the location of trainable non-zero elements in $\bm{R}$, not their actual values.

A straightforward fully connected scheme between two hierarchical layers involves $\mathcal{O}(d^2)$ connections, represented by a single dense orthogonal matrix, which leads to $d^2 - d$ edges (for instance, with $d=4$, this equates to 12 edges). As depicted in Figure~\ref{fig:info_trans}, an example of a valid matrix factorization utilizes just 10 edges, thereby requiring fewer edges than a fully connected orthogonal setup. This construction offers a useful graphical method for analyzing the parameter efficiency of OFT and aids in readily devising efficient factorizations. Notably, we draw upon the butterfly graphs~\cite{cooley1965algorithm}, a well-known topology in the Cooley-Tukey algorithm, which connects $d$ source nodes to $d$ target nodes with high density using only $\mathcal{O}(d\log d)$ edges. For comparison, while the arrangement in Figure~\ref{fig:info_trans} employs 10 edges to accomplish dense connections, the butterfly configuration achieves this with merely 8 edges. In the subsequent discussion, we elaborate on the advantages of butterfly structures for parameter-efficient modeling.

\section{Orthogonal Parameterization by Butterfly Factorization}

The butterfly network originated from the Cooley-Tukey algorithm for the efficient computation of the fast Fourier transform. In the context of Fourier transforms, modifying a frequency-domain component impacts the entire spatial domain, mirroring our challenge of disseminating information from every node in one layer to all nodes in another. Moreover, butterfly architectures are also prominent as network topologies~\cite{solihin2015fundamentals,li2011linear}, facilitating scalable and efficient communication. Considering $k\geq 2$ as a power of $2$, we define the \emph{butterfly factor} $\bm{B}^F(k)$ as

\begin{equation}\label{eq:but_factor}
\begin{aligned}
    \bm{B}^F(k)=\begin{bmatrix}
    \text{diag}(\bm{d}_1) & \text{diag}(\bm{d}_2) \\
    \text{diag}(\bm{d}_3) & \text{diag}(\bm{d}_4)
    \end{bmatrix}\in\mathbb{R}^{k\times k},
\end{aligned}
\end{equation}

with $\bm{d}_i\in\mathbb{R}^{\frac{k}{2}}$ for each $i$. For $d=2^N$, the $d$-dimensional \emph{butterfly matrix} $\bm{B}(d)\in\mathbb{R}^{d\times d}$ is defined recursively as follows:

\begin{equation}
\begin{aligned}
    \!\!\bm{B}(d)=\tilde{\bm{B}}(d,d)\!\cdot\!\begin{bmatrix}
    \bm{B}_1(\frac{d}{2})\!\!\!\!\! & \bm{0} \\
    \bm{0}\!\!\!\!\! &\bm{B}_2(\frac{d}{2})
    \end{bmatrix}= \tilde{\bm{B}}(d,d)\tilde{\bm{B}}(d,\frac{d}{2})\cdots\tilde{\bm{B}}(d,2),
\end{aligned}
\end{equation}

Here, $\bm{B}_1(\frac{d}{2})$ and $\bm{B}_2(\frac{d}{2})$ denote two butterfly matrices, each of dimension $\frac{d}{2}$. We introduce the concept of the \emph{butterfly component}, denoted as $\thickmuskip=2mu \medmuskip=2mu \tilde{\bm{B}}(d,k)=\text{diag}(\bm{B}^F_1(k),\cdots,\bm{B}^F_{d/k}(k))$, representing a block-diagonal $d\times d$ matrix where each block is of size $k$. Here, every $\bm{B}^F_i(k)$ corresponds to a butterfly factor as defined in Equation~\ref{eq:but_factor}. This formulation sets the stage for utilizing butterfly matrices to represent orthogonal matrices. To accomplish this, it is necessary to ensure that all constituent factors $\tilde{\bm{B}}(d,k)$ of the butterfly matrix $\bm{B}(d)$ are themselves orthogonal. As a starting point, we examine the case where the block-diagonal matrix $\tilde{\bm{B}}(d,2)$ consists of $2\times2$ blocks. Orthogonality for $\tilde{\bm{B}}(d,2)$ can be readily enforced by employing the Cayley transform, or equivalently, parameterizing each block as a $2$-dimensional rotation, following the methodology outlined in \cite{liu2021orthogonal,qiu2023controlling}. Since the arrangement of non-zero elements in any butterfly component corresponds to a permutation of those in $\tilde{\bm{B}}(d,2)$, all such butterfly components can be similarly parameterized as orthogonal. Consequently, this approach yields an efficient means to parametrize orthogonal matrices through a composition of several $2\times 2$ orthogonal transformations. Building on \cite{chen2022pixelated}, we extend the butterfly matrices to define a block butterfly component $\tilde{\bm{B}}^b(d,k)$, where each element in $\bm{d}_i$ is a $b\times b$ submatrix. To maintain the orthogonality of the block butterfly component $\tilde{\bm{B}}^b(d,2)$, we ensure each $2b\times 2b$ block is orthogonally parameterized. For butterfly components with $k>2$, denoted as $\tilde{\bm{B}}^b(d,k)$, the nonzero structure is simply a block-wise permutation of that for $\tilde{\bm{B}}^b(d,2)$, and therefore, these can be orthogonalized using the same procedure. Collectively, these results allow us to define the forward computation in BOFT as follows:

\begin{equation}
    \begin{aligned}\nonumber
    \bm{z}=\big(\bm{R}(m,b)\cdot\bm{W}^0\big)^\top\bm{x},~~~~\text{s.t.}~\bigg{\{}\bm{R}(m,b)=\prod_{i=1}^m \tilde{\bm{B}}^b_{(d,i)}~~\&~~\underbrace{\big(\tilde{\bm{B}}^b_{(d,j)}\big)^\top\tilde{\bm{B}}^b_{(d,j)}=\tilde{\bm{B}}^b_{(d,j)}\big(\tilde{\bm{B}}^b_{(d,j)}\big)^\top\!=\bm{I}_d}_{\forall j\in [1, m]}\bigg{\}},
    \end{aligned}
\end{equation}

In this formulation, for ease of notation, $\tilde{\bm{B}}^b(d,2^{m - i+1})$ is abbreviated as $\tilde{\bm{B}}^b_{(d,i)}$, and $\bm{I}_d$ indicates a $d \times d$ identity matrix. The orthogonal matrix $\bm{R}(m,b)\in\mathbb{R}^{d\times d}$ is constructed by multiplying several orthogonal butterfly matrices. We refer to BOFT with $\bm{R}(m,\frac{b}{2})$ as BOFT($m$,$b$), requiring $b\geq 2$. In the specific case where $m=1$, BOFT($1$,$b$) simplifies to the block-diagonal OFT~\cite{qiu2023controlling} with block size $b$. When $b=d$, BOFT($1$,$d$) recovers the standard OFT~\cite{qiu2023controlling}, yielding an unconstrained orthogonal matrix. Taking $m=\log \frac{2d}{b}$ and block butterfly matrix $\bm{B}^{\frac{b}{2}}(d)$ for $\bm{R}$ produces a dense orthogonal matrix. Generally, BOFT($m$,$b$) provides $\frac{1}{2}(b-1)dm$ meaningful parameters to adapt a linear transformation of shape $d\times n$. If the full butterfly structure is employed ($m=\log d,b=2$), the number of parameters is $\mathcal{O}(d\log d)$. Conversely, the traditional OFT with a dense orthogonal matrix utilizes $\mathcal{O}(d^2)$ parameters, while block-diagonal OFT with $r$ blocks uses $\mathcal{O}(bd)$. Notably, achieving a dense orthogonal matrix in the original OFT necessitates setting $b=d$, whereas BOFT is flexible with respect to the choice of $b$ to attain a dense transformation.

\textbf{Identity Initialization for BOFT}. A common strategy for finetuning involves initializing from the pretrained weights so that the updated model remains close to its original state. For instance, LoRA sets its low-rank parameters to zero initially. In BOFT, each butterfly factor is initialized as the identity, which, for Cayley parameterization, equates to using a zero skew-symmetric matrix.

\textbf{Multiplicative Dropout in BOFT}. In analogy to LoRA~\cite{hulora2022}, where dropout regularization is applied to the low-rank update, BOFT requires a tailored dropout approach due to its multiplicative parameterization. The additive form of standard Dropout~\cite{srivastava2014dropout} does not naturally extend to the multiplicative update in BOFT. To overcome this, we introduce a multiplicative dropout scheme tailored for BOFT: given that $\bm{R}(m,b)$ is an ordered product of $m$ butterfly matrices, and that each can be viewed as a $2b\times 2b$ block-diagonal orthogonal matrix (possibly permuted), our dropout process selects a proportion $p_1$ of the butterfly factors at random, and within each chosen butterfly, randomly sets a proportion $p_2$ of the diagonal blocks to be identity matrices, thereby regularizing the parameter updates.

\section{Intriguing Insights and Discussions}

\textbf{Expressivity of BOFT}. The butterfly architecture, when combined with suitable permutations, is capable of exactly reconstructing numerous classical fast linear transforms~\cite{dao2019learning,dao2019kaleidoscope} (\emph{e.g.}, the fast Fourier transform or the Hadamard transform). Nevertheless, it is not fully understood how closely our orthogonal butterfly matrix can approximate an arbitrary orthogonal matrix. To investigate this, we simulate the approximation of a randomly sampled dense orthogonal matrix~\cite{anderson1987generation} with dimensions $512\times 512$. The outcomes, presented in Figure~\ref{fig:para_eff}, are averaged over 10 independent runs. In the figure, the y-axis measures the approximation error, while the x-axis represents the effective count of trainable parameters. Each colored curve corresponds to BOFT with a specific block size, with the curve’s leftmost point depicting the error for BOFT($1$,$b$) (\emph{i.e.}, the original block-diagonal OFT using block size $b$). Overall, BOFT achieves superior parameter efficiency compared to OFT. For instance, BOFT($9$,$2$) attains higher expressivity than BOFT($1$,$16$) while utilizing significantly fewer parameters. Typically, choosing a smaller $b$ and larger $m$ for BOFT improves parameter efficiency. As another example, BOFT($6$,$4$) requires fewer parameters, yet approximates with an error close to that of BOFT($2$,$16$). More broadly, the butterfly matrix imposes a more structured form on the orthogonal group than the block-diagonal pattern, resulting in BOFT being demonstrably more expressive than OFT for a given block size.

\begin{theorem}[Expressivity of BOFT]\label{thm:exp_boft}
    BOFT possesses greater expressivity than OFT for identical block sizes. To achieve universal approximation over orthogonal matrices of size $d$, one can construct products of butterfly matrices in the form $\bm{B}_{d-1,1}(d)\bm{B}^\top_{d-1,2}(d)\cdots\bm{B}_{1,1}(d)\bm{B}^\top_{1,2}(d)$, with $\bm{B}_{i,j}(d), \forall i,\forall j$ being butterfly matrices.
\end{theorem}

The implication of Theorem~\ref{thm:exp_boft} is a straightforward extension for BOFT: the orthogonal matrix can be generalized as $\thickmuskip=2mu \medmuskip=2mu \bm{R}^G(m_1,b_1,m_2,b_2,l)=\bm{R}_{l,1}(m_1,b_1)\bm{R}_{l,2}^T(m_2,b_2)\cdots\bm{R}_{1,1}(m_1,b_1)\bm{R}_{1,2}^T(m_2,b_2)$, where $\bm{R}_{i,j}^T(m,b)$ denotes the orthogonal matrix in BOFT. By setting $m_1=m_2=\log d$, $b_1=b_2=2$, and $l=d-1$, this generalized matrix $\bm{R}^G(m_1,b_1,m_2,b_2,l)$ is capable of spanning the full orthogonal group. This structure is known as the kaleidoscope hierarchy~\cite{dao2019kaleidoscope}. However, it is important to observe that enhanced expressivity does not always result in improved finetuning outcomes; for example, even full finetuning, which is maximally expressive, may often yield inferior results. The crucial aspect for model generalizability is balancing expressivity against regularization. BOFT expands the parameter space for finetuning while embedding structural priors, helping to achieve a better compromise between expressivity and regularity.

\textbf{Spectral properties}. Compared to LoRA, orthogonal finetuning more effectively maintains desirable spectral properties because it exactly retains the spectral norm of the original pretrained weight matrix $\bm{W}^0$. This can be demonstrated through its singular value decomposition: $\bm{W}^0 = \bm{U}\bm{\Sigma}\bm{V}^\top$, with $\bm{U}$ and $\bm{V}$ being orthogonal matrices and $\bm{\Sigma}$ a diagonal matrix of singular values. In both OFT and BOFT, an orthogonal matrix $\bm{R}$ is applied to the left of this decomposition, producing updated weights of the form $\bm{R}\bm{U}\bm{\Sigma}\bm{V}^\top$. This operation leaves the maximal singular value—and thus the spectral norm of $\bm{W}^0$—unaltered. Maintaining this spectral norm has been demonstrated to promote both stable training and improved generalization performance~\cite{miyato2018spectral,yoshida2017spectral}. Additional mathematical aspects are discussed in Appendix~\ref{app:sn_property}.

\textbf{Orthogonal finetuning as learning bilinear similarity}. The BOFT method can be interpreted as optimizing a bilinear similarity function, specifically $\bm{w}^0_i\bm{R}\bm{x}$, where $\bm{w}^0_i$ denotes the $i$-th column (neuron) of $\bm{W}^0$. From this perspective, BOFT is equivalent to learning the matrix $\bm{R}$ in a bilinear similarity, subject to a strict orthogonality constraint, directly relating it to both distance metric learning~\cite{xing2002distance} and classical bilinear forms~\cite{roman2005advanced}.

\textbf{Inductive bias and generalization in BOFT}. In BOFT, the matrix $\bm{R}(m,b)$ typically characterizes a specific structured subset within the orthogonal group, thereby constraining the space of learnable hypotheses and introducing a distinct inductive bias. We posit that the structured inductive bias imparted by butterfly factorization enhances generalization capabilities, given that this structure underlies a number of foundational linear transformations—such as the discrete Fourier transform, discrete sine/cosine transform, and Hadamard transform—recognized for their effectiveness~\cite{dao2019kaleidoscope}. Furthermore, the use of sparse matrix factorizations in BOFT may introduce additional, possibly implicit, forms of inductive bias~\cite{gunasekar2017implicit,li2018algorithmic,liu2019neural}.

\textbf{Relation to butterfly-based sparse training}. Numerous studies~\cite{dao2019learning,dao2019kaleidoscope,chen2022pixelated,dao2022monarch} have explored sparse training using the butterfly parameterization. These approaches generally involve directly reparameterizing the weight matrices via the butterfly parameterization and then training neural architectures from the beginning. In \cite{dao2022monarch}, the method involves adapting pretrained weights by projecting them onto a variation of butterfly matrices, followed by finetuning these projected elements for specific downstream tasks. In contrast, BOFT introduces a distinct finetuning mechanism, applying layer-shared weight matrices to transform the weights during adaptation.

\input{rebuttal-results}

\section{Concluding Remarks and Limitations}

This work presents Orthogonal Butterfly, a broadly applicable, parameter-efficient finetuning technique for foundation models, constructed upon the butterfly structure. The principal idea underpinning its improved parameter efficiency is to represent a dense orthogonal matrix as a product of several sparse orthogonal matrices. To facilitate the discovery of suitable matrix decompositions, we introduce a graphical information transmission perspective. Within this context, the butterfly structure emerges as a powerful tool for fulfilling our requirements in sparse orthogonal matrix factorizations. We empirically establish the utility of BOFT for finetuning a range of model types, including large language models, expansive vision architectures, and text-to-image generative systems. Our experimental results further affirm BOFT’s advantage as a universal approach for model finetuning.

Nonetheless, BOFT does present certain limitations. Because the ultimate orthogonal matrix in BOFT results from multiplying several orthogonal matrices, its training incurs a slightly greater computational cost relative to OFT. Finding methods to reduce the training overhead for BOFT constitutes an unresolved challenge. Importantly, after finetuning is completed, the orthogonal matrices generated by BOFT may be incorporated directly into the pretrained model without introducing any additional delay at inference time. Additionally, it remains an open question whether the butterfly topology represents the most information-efficient structure for transmission. Our proposed information transmission framework bridges to another line of research—namely, computer networking—where the design and efficiency of network topologies for data transmission are well-explored. This connection paves the way for potentially adopting more efficient network configurations to construct dense orthogonal matrices.